\subsection{紧空间的逆极限}

命题 8. 设 \((\mathbf{X}_{\alpha}, f_{\alpha \beta})\) 是由有向集 I 加标的紧空间的一个逆系统，使得 \(f_{\alpha \alpha}\) 对每个 \(\alpha \in \mathbf{I}\) 是恒同映射。设 \(\mathbf{X} = \varprojlim \mathbf{X}_{\alpha}\) 是逆极限，\(f_{\alpha}: \mathbf{X} \to \mathbf{X}_{\alpha}\) 是典范映射（§ 4，第 4 目）。则

\begin{enumerate}
\def\labelenumi{\alph{enumi})}
\tightlist
\item
  X 是紧的，且对每个 \(\alpha \in I\) 我们有
\end{enumerate}

\[
f _ {\alpha} (\mathbf {X}) = \bigcap_ {\beta \geqslant \alpha} f _ {\alpha \beta} (\mathbf {X} _ {\beta}).\tag{I}
\]

\begin{enumerate}
\def\labelenumi{\alph{enumi})}
\setcounter{enumi}{1}
\tightlist
\item
  若诸 \(\mathbf{X}_{\alpha}\) 都非空，则 \(\mathbf{X}\) 非空。
\end{enumerate}

X 是 \(\prod_{\alpha} X_{\alpha}\) 的一个闭子空间（§ 8，第 2 目，命题 7，推论 2）

而后者由第 5 目定理 3 与第 3 目命题 3 是紧的。其余断言是《集合论》第三章 § 7 第 4 目定理 1 的推论。我们应用这个定理时，取 \(\mathfrak{S}_{\alpha}\) 为 \(\mathbf{X}_{\alpha}\) 的闭子集组成的集合。条件 (i) 与 (ii) 分别就是公理 \((\mathrm{O}_{\mathbf{i}}^{\prime})\) 与 \((\mathbf{C}^{\prime \prime})\)；条件 (iii) 得以满足，因为 \(\{x_{\alpha}\}\) 是闭的且 \(f_{\alpha \beta}\) 连续（\(\S 2\)，第 1 目，定理 1）；最后，条件 (iv) 由第 4 目定理 2 的推论 2 得以满足。

推论 1. 设 \((\mathbf{X}_{\alpha}, f_{\alpha\beta})\) 是由一个有向集加标的拓扑空间的逆系统，使得对每对指标 \(\alpha\), \(\beta\)（满足 \(\alpha \leqslant \beta\)）及每个 \(x_{\alpha} \in X_{\alpha}\)，\(\overline{f}_{\alpha\beta}(x_{\alpha})\) 是紧的。则等式 (1) 成立，且 \(\overline{f}_{\alpha}(x_{\alpha})\) 对每个 \(\alpha\) 与每个 \(x_{\alpha} \in X_{\alpha}\) 是紧的。

对每个 \(x_{\alpha} \in \bigcap_{\beta \geqslant \alpha} f_{\alpha \beta}(X_{\beta})\) 及每个 \(\beta \geqslant \alpha\)，令 \(L_{\beta}\) 表示 \(\overline{f}_{\alpha \beta}^{1}(x_{\alpha})\)。若 \(\alpha \leqslant \beta \leqslant \gamma\)，则我们有 \(f_{\beta \gamma}(L_{\gamma}) \subset L_{\beta}\)，且满足 \(\beta \geqslant \alpha\) 的全体指标组成的集合在指标集中是共尾的。由此立即得出，诸 \(L_{\beta} (\beta \geqslant \alpha)\) 构成拓扑空间的一个逆系统（以诸 \(f_{\beta \gamma}\) 的限制为映射），其逆极限 \(L\) 同胚于 \(\overline{f}_{\alpha}^{1}(x_{\alpha})\)。既然由假设诸 \(L_{\beta}\) 是紧的且非空，该推论由命题 8 得出。

推论 2. 设 \((\mathbf{X}_{\alpha}, f_{\alpha \beta})\) 与 \((\mathbf{X}_{\alpha}^{\prime}, f_{\alpha \beta}^{\prime})\) 是由同一有向集 I 加标的拓扑空间的两个逆系统，并设 \((u_{\alpha})\) 是映射 \(u_{\alpha}: \mathbf{X}_{\alpha} \to \mathbf{X}_{\alpha}^{\prime}\) 的一个逆系统。设 \(\mathbf{X} = \varprojlim \mathbf{X}_{\alpha}\)，\(\mathbf{X}^{\prime} = \varprojlim \mathbf{X}_{\alpha}^{\prime}\)，\(u = \varprojlim u_{\alpha}\)。则：

\begin{enumerate}
\def\labelenumi{\alph{enumi})}
\item
  若 \(x' = (x_{\alpha}') \in \mathbf{X}'\) 使 \(\overline{u}_{\alpha}(x_{\alpha}')\) 对每个 \(\alpha \in I\) 是紧的且非空，则 \(\overline{u}^1(x')\) 是紧的且非空。
\item
  若诸 \(X_{\alpha}\) 是紧的，诸 \(X_{\alpha}^{\prime}\) 是豪斯多夫的，且诸 \(u_{\alpha}\) 是满射与连续的，则 u 是满射。
\end{enumerate}

令 \(\mathbf{L}_{\alpha}\) 表示 \(\overline{\boldsymbol{u}}_{\alpha}^{1}(x_{\alpha}^{\prime})\)；则显然诸 \(\mathbf{L}_{\alpha}\) 构成拓扑空间的一个逆系统（以诸 \(f_{\alpha \beta}\) 的限制为映射），且 \(\overline{\boldsymbol{u}}^{1}(x^{\prime} = \mathbf{L})\) 是诸 \(\mathbf{L}_{\alpha}\) 的逆极限；因此断言 a) 由命题 8 得出。鉴于第 3 目命题 3，断言 b) 是一个直接推论。

